%==============================================================================
%  CHAPITRE 5 - Plan de releases
%==============================================================================
\chapter{Plan de releases}
\label{chap:releases}

Le versionnage suit \textbf{Semantic Versioning} (SemVer :
MAJOR.MINOR.PATCH). Quatre releases jalonnent le projet, une à l'issue de
chaque paire de sprints.

\section{Synthèse}

\begin{center}
\begin{tabular}{lllll}
\toprule
\textbf{Version} & \textbf{Nom} & \textbf{Date} & \textbf{Sprints} & \textbf{Statut} \\
\midrule
v0.1.0 & MVP --- Fondation          & 25/08/2026 & S1, S2 & Livrée \\
v0.2.0 & Visites \& collaboration   & 22/09/2026 & S3, S4 & Livrée \\
v0.3.0 & Analytics \& dashboards    & 20/10/2026 & S5, S6 & Livrée \\
v1.0.0 & Release Candidate          & 16/11/2026 & S7, S8 & Livrée\footnotemark \\
v1.1.0 & Exploitation \& spatial    & 18/01/2027 & S9--S11 & Livrée \\
v1.2.0 & Sécurité \& PWA livrable   & 15/02/2027 & S12, S13 & Livrée \\
v1.3.0 & Conformité miel            & 15/03/2027 & S14, S15 & Livrée \\
v1.4.0 & Sessions serveur \& portée & 29/03/2027 & S16 & Livrée \\
v1.5.0 & Dettes soldées             & 26/04/2027 & S17, S18 & Livrée \\
\bottomrule
\end{tabular}
\footnotetext{Sous réserve : les livrables documentaires US-039 et US-040
restent à produire. Voir l'encadré en fin de chapitre.}
\end{center}

\section{v0.1.0 --- MVP : Fondation}

\begin{center}
\begin{tabular}{lp{11cm}}
\toprule
Contenu &
  CRUD complet des entités métier ;
  architecture Spring Boot multi-niveaux ;
  composition de la ruche (corps/hausse/cadre) ;
  configuration \texttt{ConfigZumm.ini} ;
  i18n FR/EN/AR ;
  sécurité TLS 1.3 / X.509 ;
  harnais de tests d'intégration Testcontainers. \\
Critère d'acceptation & Application déployable en HTTPS avec données de test. \\
Statut & 	extbf{Livrée}. \\
\bottomrule
\end{tabular}
\end{center}

\section{v0.2.0 --- Visites \& collaboration}

\begin{center}
\begin{tabular}{lp{11cm}}
\toprule
Contenu &
  Workflow planification/approbation de visite ;
  rapport de visite complet avec photos ;
  mode hors-ligne PWA ;
  authentification OAuth 2.0 et repli local ;
  contrôle d'accès RBAC complet. \\
Critère d'acceptation & Parcours utilisateur complet en ligne et hors-ligne. \\
Statut & 	extbf{Livrée}. \\
\bottomrule
\end{tabular}
\end{center}

\section{v0.3.0 --- Analytics \& dashboards}

\begin{center}
\begin{tabular}{lp{11cm}}
\toprule
Contenu &
  Calendrier agents $\times$ ruches ;
  tableaux de bord production \& alertes ;
  synthèse ROI ;
  rappels et export CSV/TXT ;
  ingestion des mesures capteurs ;
  contexte météo local ;
  service web API (REST/OpenAPI). \\
Critère d'acceptation & Dashboards fonctionnels avec données réelles. \\
Statut & 	extbf{Livrée}. \\
\bottomrule
\end{tabular}
\end{center}

\section{v1.0.0 --- Release Candidate}

\begin{center}
\begin{tabular}{lp{11cm}}
\toprule
Contenu &
  Fonctionnalités avancées (carte et rayons de butinage, suivi de reine,
  QR code de traçabilité) ;
  détection d'anomalie adaptative (IA) ;
  tests complets (unitaires/intégration/charge) ;
  documentation \& soutenance ;
  monitoring \& alerting. \\
Critère d'acceptation & Solution complète prête pour la production. \\
Statut & 	extbf{Livrée sous réserve} --- voir l'encadré ci-dessous. \\
\bottomrule
\end{tabular}
\end{center}

\begin{encadre}
\textbf{Réserve sur la v1.0.0.} Le périmètre \emph{applicatif} est livré,
mais la ligne « documentation \& soutenance » ne l'est pas : US-039
(diagrammes UML) et US-040 (rapport, poster, présentation) restent à produire,
la charte académique de l'épreuve interdisant de les générer. \textbf{21 des
37 points du sprint 8 ne sont donc pas acquis}, et la version ne peut pas être
déclarée livrée sans cette réserve.
\end{encadre}

\section{Après la v1.0.0 --- cinq versions de consolidation}

Les sprints 9 à 18 n'ajoutent pas de version majeure : ils durcissent et
soldent un produit déjà livrable. Ils se répartissent en cinq versions
mineures, toutes \textbf{livrées}.

\begin{center}
\begin{tabular}{llp{7.2cm}}
\toprule
\textbf{Version} & \textbf{Sprints} & \textbf{Contenu} \\
\midrule
v1.1.0 & S9--S11  & Notifications, rapport PDF, audit, prévision ; requêtes spatiales PostGIS ; socle de test du front \\
v1.2.0 & S12--S13 & Durcissement de la sécurité (29 vulnérabilités $\rightarrow$ 0) ; PWA déployée, graphiques et carte \\
v1.3.0 & S14--S15 & Conformité miel (directive UE 2024/1438), idempotence ; ressources de langue externalisées \\
v1.4.0 & S16      & BFF : plus aucun jeton dans le navigateur (ADR-006) ; autorisation horizontale (US-057) \\
v1.5.0 & S17--S18 & Dettes techniques soldées ; agrégats calculés en base ; parité de types client/contrat \\
\bottomrule
\end{tabular}
\end{center}

\begin{encadre}
\textbf{Règle de livraison.} Une release n'est publiée que si tous les
critères de la Definition of Done (\S1.2.4) sont satisfaits pour l'ensemble
des stories concernées, et si le pipeline CI/CD (chapitre~\ref{chap:devops})
est intégralement vert sur la branche \texttt{release/*} correspondante.
\end{encadre}
